% 型9：難易度別総合演習型
% 難易度記号と目標時間を付けた一覧表を先頭に置き，読者に選ばせる．
% 単元を横断する総復習に向く．既存の \MMLegend と \MMDifficulty を使う．
\documentclass[lualatex,paper=b5j,fontsize=10pt,twoside]{jlreq}

\usepackage{surikumi-mondaishu}

\MathMagazineSetup{
  feature={},
  series={総合演習／数学 I・A},
  series-position=left,
  pattern=diamonds,
  title={分野をまたぐ 6 題},
  author={入試基礎講座},
  author-font=mincho,
  author-position=right,
  title-fill=black,
  deck={難易度と目標時間を見て，取り組む問題を自分で選ぶ．全問を順に解く必要はない．}
}

\begin{document}

\MakeMathMagazineTitle

\MMLegend{A は基本，B は標準，C はやや難を表す．
$*$ は 10 分，\MMFiveMinuteMark{} は 5 分を表し，
$**$\MMFiveMinuteMark{} なら目標 25 分である．}

\begin{msindex}
  1 & 二次不等式と整数解 & \MMDifficulty{A}{10} & 基本 \\
  2 & 三角比と面積 & \MMDifficulty{A}{15} & 09 東京電機大・改 \\
  3 & 場合の数（重複を除く） & \MMDifficulty{B}{20} & 12 立命館大・改 \\
  4 & データの分析と分散 & \MMDifficulty{B}{20} & 基本 \\
  5 & 整数の性質（合同式） & \MMDifficulty{B}{25} & 14 千葉大・改 \\
  6 & 図形と論証 & \MMDifficulty{C}{30} & 08 一橋大・改 \\
\end{msindex}

\begin{mmproblems}
  \item 二次不等式 $x^2-(a+1)x+a<0$ を満たす整数 $x$ がちょうど 2 個
    存在するような定数 $a$ の範囲を求めよ．
    \MMSource{\MMDifficulty{A}{10}}

  \item $\triangle\mathrm{ABC}$ において $\mathrm{AB}=5$，$\mathrm{BC}=7$，
    $\mathrm{CA}=8$ である．$\triangle\mathrm{ABC}$ の面積を求めよ．
    \MMSource{\MMDifficulty{A}{15}\quad 09 東京電機大・改}

  \item 赤，青，白の玉が各 3 個ずつある．同じ色の玉は区別しない．
    この 9 個の玉から 4 個を選ぶ方法は何通りあるか．
    \MMSource{\MMDifficulty{B}{20}\quad 12 立命館大・改}

  \item 5 個の値 $2,\,4,\,6,\,8,\,10$ の分散を求めよ．
    また，各値を 3 倍したときの分散を求めよ．
    \MMSource{\MMDifficulty{B}{20}}

  \item $7^{100}$ を $10$ で割った余りを求めよ．
    \MMSource{\MMDifficulty{B}{25}\quad 14 千葉大・改}

  \item 鋭角三角形 $\mathrm{ABC}$ の内部に点 $\mathrm{P}$ をとる．
    $\mathrm{PA}+\mathrm{PB}+\mathrm{PC}$ が最小となる $\mathrm{P}$ について，
    $\angle\mathrm{APB}=\angle\mathrm{BPC}=\angle\mathrm{CPA}$ が成り立つことを示せ．
    \MMSource{\MMDifficulty{C}{30}\quad 08 一橋大・改}
\end{mmproblems}

\MMHeading{略\qquad 解}

\begin{mmproblems}
  \item $x^2-(a+1)x+a=(x-1)(x-a)$ である．$a>1$ のとき解は $1<x<a$ で，
    整数解が $2,\,3$ の 2 個となるのは $3<a\leq 4$ のときである．
    $a<1$ のとき解は $a<x<1$ で，整数解が $-1,\,0$ の 2 個となるのは
    $-2\leq a<-1$ のときである．よって
    \[
      -2\leq a<-1,\quad 3<a\leq 4
    \]

  \item 余弦定理より $\cos\mathrm{B}=\dfrac{25+49-64}{2\cdot 5\cdot 7}=\dfrac{1}{7}$
    である．よって $\sin\mathrm{B}=\dfrac{4\sqrt{3}}{7}$ となり，面積は
    \[
      \dfrac{1}{2}\cdot 5\cdot 7\cdot\dfrac{4\sqrt{3}}{7}=10\sqrt{3}
    \]

  \item 各色の個数を $(x,y,z)$ とすると $x+y+z=4$，$0\leq x,y,z\leq 3$ である．
    制限なしの解は $15$ 通りで，いずれかが $4$ になる場合が $3$ 通りあるから
    \[
      15-3=12
    \]

  \item 平均は $6$ で，分散は $\dfrac{16+4+0+4+16}{5}=8$ である．
    各値を $3$ 倍すると分散は $3^2$ 倍になるから
    \[
      72
    \]

  \item $7^1,7^2,7^3,7^4$ の一の位は $7,9,3,1$ と 4 個周期で繰り返す．
    $100=4\cdot 25$ であるから
    \[
      1
    \]

  \item 点 $\mathrm{P}$ を固定して考えると，$\mathrm{PA}+\mathrm{PB}+\mathrm{PC}$ が
    最小となる点はフェルマー点である．$\triangle\mathrm{APB}$ を $\mathrm{P}$ の
    まわりに $60^\circ$ 回転して折れ線を一直線に伸ばす議論により，
    三つの角が等しくなることが従う．
\end{mmproblems}

\end{document}
